%% ch14.tex --- Anomaly Detection
%% Numbers are computed by verify/ch14_anomaly.py.
%% E2 (engines) and DC (data-centre machines) are shared with the companion
%% lecture.  Figure coordinates are generated from the same script's data.

\chapter{Anomaly Detection}
\label{ch:ch14}

\epigraph{Normal is an illusion. What is normal for the spider is chaos for the
fly.}{---\ attributed to Morticia Addams}

\begin{chapterlearning}
After working through this chapter you will be able to:
\begin{enumerate}
  \item Frame anomaly detection as density estimation, and say where it is used.
  \item State the Gaussian density and derive the maximum-likelihood estimates
        of its mean and variance.
  \item Build an anomaly detector from one Gaussian per feature, and flag
        examples whose density falls below a threshold $\varepsilon$.
  \item Evaluate the detector on a small labelled validation set and choose
        $\varepsilon$ by $F_1$, accounting for the skew of \cref{ch:ch11}.
  \item Decide between anomaly detection and supervised learning for a given
        problem.
  \item Transform features towards a Gaussian shape, and invent features from the
        anomalies the detector misses.
  \item Use the multivariate Gaussian, and explain how it relates to the
        independent model.
\end{enumerate}
\end{chapterlearning}

\Cref{ch:ch13} asked what the typical examples look like and grouped them.
Anomaly detection asks the opposite question: which examples look \emph{unlike}
the rest? A manufacturer wants to catch the faulty engine before it ships; a bank
the fraudulent transaction before it clears; an operations team the failing
server before it falls over. In each case almost all the data is normal, the
anomalies are rare and varied, and many of them have never been seen before.

The approach of this chapter is to learn what normal looks like, as a probability
density, from unlabelled data that is mostly normal --- and then to flag anything
that the density says is improbable.

\section{Finding Unusual Events}
\label{sec:unusual}

\begin{defbox}{Dataset E2}
Aircraft engines described by two features measured at the end of production:
$x_1$, the heat generated ($^\circ$C), and $x_2$, the vibration intensity
(mm/s). There are $6000$ engines that passed inspection, used for training, and
validation and test sets of $2000$ good and $10$ faulty engines each. The data is
synthetic, generated so that the truth is known: good engines are Gaussian
around $(500, 4)$, and faulty ones lie between $3.2$ and $5$ standard deviations
from the centre, in random directions.
\end{defbox}

\begin{defbox}{Anomaly detection by density estimation}
Given unlabelled examples $\xii{1},\dots,\xii{m}$, assumed to be normal, build a
model $p(\vx)$ of the probability density of the features. For a new example
$\vx_{\text{test}}$,
\[
  p(\vx_{\text{test}}) < \varepsilon \;\Longrightarrow\; \text{flag as an
  \term{anomaly}},
\]
for a threshold $\varepsilon$ chosen as in \cref{sec:ad-eval}.
\end{defbox}

The same recipe serves many settings. In fraud detection the features describe a
user's activity --- logins per day, pages visited, typing speed, purchase
amounts --- and an account whose behaviour is improbable is sent for review. In
manufacturing they are measurements of each unit off the line. In monitoring
they are a machine's memory use, disk accesses and CPU load.

\section{The Gaussian Distribution}
\label{sec:gaussian}

\begin{defbox}{Gaussian (normal) distribution}
A real random variable $x$ is \term{Gaussian} with mean $\mu$ and variance
$\sigma^2$, written $x \sim \mathcal{N}(\mu, \sigma^2)$, if its density is
\[
  p(x;\mu,\sigma^2) = \frac{1}{\sqrt{2\pi}\,\sigma}
                      \exp\!\Bigl(-\frac{(x-\mu)^2}{2\sigma^2}\Bigr).
\]
$\mu$ is the centre of the bell and $\sigma$ its width. The area under the curve
is $1$.
\end{defbox}

\begin{figure}[htbp]
  \centering
  \begin{tikzpicture}
    \begin{axis}[mlaxis, width=0.74\linewidth, height=4.8cm,
      xlabel={$x$}, ylabel={$p(x)$}, xmin=-4, xmax=6, ymin=0, ymax=0.85,
      legend style={font=\scriptsize, at={(0.98,0.97)}, anchor=north east}]
      \addplot[MLBlue, very thick, domain=-4:6, samples=120] {exp(-x^2/2)/sqrt(2*pi)};
      \addlegendentry{$\mu=0,\ \sigma=1$}
      \addplot[MLRose, very thick, domain=-4:6, samples=120] {exp(-x^2/0.5)/sqrt(0.5*pi)};
      \addlegendentry{$\mu=0,\ \sigma=0.5$}
      \addplot[MLTeal, very thick, domain=-4:6, samples=120] {exp(-x^2/8)/sqrt(8*pi)};
      \addlegendentry{$\mu=0,\ \sigma=2$}
      \addplot[MLAmber, thick, dashed, domain=-4:6, samples=120] {exp(-(x-3)^2/0.5)/sqrt(0.5*pi)};
      \addlegendentry{$\mu=3,\ \sigma=0.5$}
    \end{axis}
  \end{tikzpicture}
  \caption{Gaussian densities. Halving $\sigma$ makes the bell half as wide and
  twice as tall, so that the area stays $1$; changing $\mu$ slides it along the
  axis.}
  \label{fig:gaussians}
\end{figure}

\begin{pitfall}{A density is not a probability}
With $\sigma = 0.5$ the density peaks at $0.798$; with $\sigma = 0.1$ it peaks at
$3.99$. A density can exceed $1$. Only areas under it are probabilities, so
$p(\vx)$ is meaningful for \emph{comparing} examples --- which is all anomaly
detection needs --- but not as ``the probability that $\vx$ is normal''.
\end{pitfall}

\begin{derivation}{Maximum-likelihood estimates of $\mu$ and $\sigma^2$}
Given independent samples $\xii{1},\dots,\xii{m}$ from $\mathcal{N}(\mu,\sigma^2)$,
the log-likelihood is
\[
  \ell(\mu,\sigma^2) = \sum_{i=1}^{m}\ln p(\xii{i};\mu,\sigma^2)
  = -\frac m2\ln(2\pi\sigma^2) - \frac{1}{2\sigma^2}\sum_{i=1}^{m}\bigl(\xii{i} - \mu\bigr)^2 .
\]
Set its derivative with respect to $\mu$ to zero:
\[
  \frac{\partial\ell}{\partial\mu} = \frac{1}{\sigma^2}\sum_{i}\bigl(\xii{i} - \mu\bigr) = 0
  \quad\Longrightarrow\quad
  \hat\mu = \frac1m\sum_{i=1}^{m}\xii{i} .
\]
Then with respect to $v = \sigma^2$:
\[
  \frac{\partial\ell}{\partial v} = -\frac{m}{2v} + \frac{1}{2v^2}\sum_i\bigl(\xii{i} - \hat\mu\bigr)^2 = 0
  \quad\Longrightarrow\quad
  \hat\sigma^2 = \frac1m\sum_{i=1}^{m}\bigl(\xii{i} - \hat\mu\bigr)^2 .
\]
Fitting a Gaussian is two averages. Statistics texts often divide by $m-1$
instead, which makes $\hat\sigma^2$ unbiased; for the thousands of examples
typical in anomaly detection the difference is negligible. On E2's heat feature
it changes $\hat\sigma^2$ from $630.96$ to $631.07$.
\end{derivation}

\section{The Anomaly Detection Algorithm}
\label{sec:ad-alg}

With $n$ features, model each with its own Gaussian and multiply.

\begin{algorithmbox}{Anomaly detection with independent Gaussians}
\begin{lstlisting}[style=pseudocode]
choose features x_1..x_n that may indicate anomalies
fit:      mu_j     = mean of x_j over the training set
          sigma2_j = mean of (x_j - mu_j)^2
score:    p(x) = product over j of N(x_j; mu_j, sigma2_j)
decide:   anomaly if p(x) < epsilon
\end{lstlisting}
Multiplying the densities treats the features as statistically independent.
The algorithm often works well even when they are not; \cref{sec:multivariate}
shows what to do when the dependence matters.
\end{algorithmbox}

\begin{workedexample}{Scoring three new engines}
From the $6000$ training engines: $\hat\mu_1 = 499.79$, $\hat\sigma_1 = 25.12$ for
heat; $\hat\mu_2 = 3.993$, $\hat\sigma_2 = 0.495$ for vibration.

\begin{center}
\small
\renewcommand{\arraystretch}{1.15}
\begin{tabular}{@{}lrrrrr@{}}
\toprule
engine & $z_1$ & $z_2$ & $p(x_1)$ & $p(x_2)$ & $p(\vx)$\\
\midrule
$(505,\ 4.1)$ & $0.21$ & $0.22$ & $0.0155$ & $0.787$ & $1.2\times10^{-2}$\\
$(560,\ 4.9)$ & $2.40$ & $1.83$ & $0.0009$ & $0.151$ & $1.4\times10^{-4}$\\
$(590,\ 5.6)$ & $3.59$ & $3.24$ & $0.000025$ & $0.0042$ & $1.0\times10^{-7}$\\
\bottomrule
\end{tabular}
\end{center}

Here $z_j = (x_j - \hat\mu_j)/\hat\sigma_j$, the number of standard deviations
from the mean. With $\varepsilon = 10^{-5}$ the first two engines are normal and
the third is flagged. A product is small when \emph{any} factor is small: one
improbable reading is enough.
\end{workedexample}

\begin{figure}[htbp]
  \centering
  \begin{tikzpicture}
    \begin{axis}[mlaxis, width=0.84\linewidth, height=6.0cm, view={0}{90},
      xlabel={heat $x_1$ ($^\circ$C)}, ylabel={vibration $x_2$ (mm/s)},
      xmin=380, xmax=620, ymin=1.8, ymax=6.2,
      legend style={font=\scriptsize, at={(0.02,0.98)}, anchor=north west}]
      \addplot[only marks, mark=*, mark size=0.8pt, MLBlue!45, forget plot] coordinates {(528.4,3.772)(509.5,3.686)(485.1,4.993)(474.0,4.142)(503.3,4.110)(514.0,4.218)(457.4,3.908)(493.9,3.909)(498.5,4.284)(457.8,4.709)(476.9,4.195)(536.4,4.242)(470.0,3.937)(468.1,4.541)(500.5,2.936)(544.8,4.672)(504.4,4.018)(521.1,4.382)(500.6,3.523)(490.2,4.044)(517.2,3.194)(495.0,4.542)(496.4,3.521)(594.1,4.166)(508.2,4.034)(483.2,3.610)(466.3,4.308)(528.3,3.986)(478.4,3.658)(527.3,5.004)(422.9,4.180)(473.8,4.050)(491.9,3.752)(468.9,4.429)(529.8,3.680)(474.5,4.020)(508.5,3.188)(514.8,3.633)(501.3,3.942)(502.1,3.299)(494.8,3.837)(499.3,4.248)(496.5,3.926)(493.8,4.584)(466.0,4.203)(539.1,4.173)(495.7,3.860)(537.8,4.476)(517.2,4.388)(526.9,4.122)(467.7,4.173)(483.9,3.011)(526.7,2.948)(506.2,4.551)(467.9,4.268)(513.2,3.650)(474.8,3.619)(532.2,4.055)(492.3,3.839)(457.2,3.994)(509.6,5.191)(489.0,3.872)(514.2,4.722)(507.8,4.591)(525.1,3.773)(489.7,4.257)(484.0,4.399)(473.4,2.920)(514.8,3.846)(541.3,4.467)(507.7,2.860)(482.3,3.513)(491.6,4.051)(474.0,4.160)(516.1,4.114)(505.6,3.007)(516.5,3.128)(514.7,3.858)(466.6,4.203)(510.1,3.698)(462.8,3.623)(457.6,3.319)(478.5,4.343)(503.2,3.406)(500.4,4.258)(465.5,4.475)(490.7,3.468)(480.3,3.971)(527.0,4.762)(522.4,3.865)(536.2,3.854)(465.2,4.115)(486.6,4.442)(507.4,4.055)(453.4,4.721)(447.1,3.562)(483.5,3.794)(469.9,3.229)(523.5,4.661)(513.4,4.585)(479.5,4.936)(493.8,4.193)(475.6,2.845)(506.9,4.392)(517.1,3.842)(487.9,4.466)(516.8,3.780)(510.8,4.337)(503.2,4.209)(523.7,4.347)(523.6,4.358)(524.2,3.480)(520.5,4.020)(470.7,3.530)(523.9,3.631)(489.7,4.210)(472.2,3.981)(551.9,3.962)(467.5,3.901)(493.4,4.581)(496.4,3.518)(456.1,3.708)(495.3,2.749)(477.2,5.048)(558.8,3.448)(495.6,4.019)(480.7,4.508)(476.9,4.722)(492.5,4.461)(529.7,3.306)(516.5,4.253)(478.8,3.789)(498.0,4.149)(483.2,4.595)(495.0,3.886)(490.5,3.448)(497.4,4.137)(503.5,3.892)(447.0,3.903)(516.8,3.594)(509.8,3.969)(496.3,2.734)(509.9,4.055)(528.5,3.708)(517.4,3.560)(465.2,4.308)(482.3,4.069)(533.8,4.242)(547.3,3.253)(465.8,3.274)(513.1,3.120)(511.5,5.008)(512.9,3.934)(452.6,3.838)(498.8,3.875)(564.0,4.533)(495.9,4.344)(553.1,4.103)(503.3,3.777)(503.1,4.300)(512.8,4.480)(526.2,4.215)(504.0,3.817)(573.2,3.479)(503.5,4.184)(533.0,3.486)(478.1,3.846)(520.6,3.827)(550.2,4.183)(515.8,4.065)(495.7,4.688)(533.4,4.037)(506.0,4.342)(510.5,3.914)(450.4,3.962)(520.8,3.813)(471.1,3.342)(515.8,3.711)(514.9,4.504)(493.8,4.051)(515.0,4.252)(543.7,3.706)(524.2,3.930)(530.3,3.755)(461.7,3.761)(487.1,2.970)(518.1,3.617)(522.4,4.190)(523.1,3.885)(507.8,4.793)(498.8,3.809)(477.8,4.299)(518.7,3.685)(504.0,4.025)(511.3,3.650)(513.6,3.999)(472.0,3.688)(501.5,4.150)(560.0,3.990)(498.1,3.802)(503.3,4.190)(505.2,3.789)(480.7,4.675)(534.2,5.154)(487.0,2.933)(458.4,3.697)(494.0,3.646)(447.4,3.076)(515.6,4.200)(505.9,3.122)(505.1,3.271)(477.9,3.307)(511.6,3.504)(545.9,3.422)(493.9,3.586)(479.8,3.639)(547.4,3.922)(498.9,4.104)(474.5,4.466)(514.9,3.367)(486.7,3.512)(497.7,3.637)(482.8,4.232)(507.7,4.213)(546.2,3.826)(522.8,3.379)(472.3,3.392)(475.3,4.088)(498.1,4.308)(483.0,3.802)(532.4,3.449)(506.0,3.023)(502.3,2.962)(498.4,3.974)(527.6,5.122)(472.3,3.052)(515.3,4.610)(500.4,5.095)(493.7,4.552)(498.7,3.688)(518.1,4.263)(509.3,3.544)(543.8,3.805)(467.4,4.204)(500.3,4.821)(454.9,3.920)(492.8,4.723)(478.0,3.726)(489.8,4.381)(483.4,4.323)};
      \costcontour{contour/draw color=MLAmber, very thick, forget plot, domain=380:620, y domain=1.8:6.2}{5.0975}{((x-499.7865)/25.1190)^2+((y-3.99277)/0.49535)^2}
      \addlegendimage{MLAmber, very thick}\addlegendentry{$\varepsilon = 10^{-3}$}
      \costcontour{contour/draw color=MLGray, very thick, forget plot, domain=380:620, y domain=1.8:6.2}{9.7027}{((x-499.7865)/25.1190)^2+((y-3.99277)/0.49535)^2}
      \addlegendimage{MLGray, very thick}\addlegendentry{$\varepsilon = 10^{-4}$}
      \costcontour{contour/draw color=MLTeal, very thick, forget plot, domain=380:620, y domain=1.8:6.2}{14.3078}{((x-499.7865)/25.1190)^2+((y-3.99277)/0.49535)^2}
      \addlegendimage{MLTeal, very thick}\addlegendentry{$\varepsilon = 10^{-5}$}
      \addplot[only marks, mark=x, mark size=2.6pt, MLRose, very thick, forget plot] coordinates {(386.4,3.831)(599.2,3.362)(558.5,5.493)(609.5,3.268)(464.5,5.733)(411.0,4.930)(540.4,2.480)(417.4,5.060)(415.1,3.453)(590.4,4.316)};
    \end{axis}
  \end{tikzpicture}
  \caption{$250$ of the $6000$ training engines (dots) and the $10$ faulty
  validation engines (crosses), with the boundaries $p(\vx) = \varepsilon$ for
  three thresholds --- ellipses aligned with the axes, drawn from the formula. A
  smaller $\varepsilon$ gives a larger ellipse: fewer false alarms and more missed
  faults. At $\varepsilon = 10^{-5}$ three faulty engines lie inside.}
  \label{fig:engines}
\end{figure}

The boundary $p(\vx) = \varepsilon$ is an ellipse because taking logarithms turns
the product into a sum: $p(\vx) = \varepsilon$ is equivalent to
$z_1^2 + z_2^2 = -2\ln\bigl(2\pi\hat\sigma_1\hat\sigma_2\,\varepsilon\bigr)$, a
circle in standardised coordinates. For $\varepsilon = 10^{-5}$ its radius is
$3.78$ standard deviations.

\section{Developing and Evaluating a Detector}
\label{sec:ad-eval}

Anomaly detection is trained without labels, but it cannot be evaluated or tuned
without some. The standard arrangement uses a few labelled anomalies, kept out of
training:

\begin{center}
\small
\renewcommand{\arraystretch}{1.15}
\begin{tabular}{@{}lrr@{}}
\toprule
set & good engines & faulty engines\\
\midrule
training   & $6000$ & $0$ (assumed)\\
validation & $2000$ & $10$\\
test       & $2000$ & $10$\\
\bottomrule
\end{tabular}
\end{center}

Fit $\mu_j$ and $\sigma_j^2$ on the training set. On the validation set, predict
$y = 1$ (anomaly) when $p(\vx) < \varepsilon$, and choose $\varepsilon$ --- and,
if you are comparing them, the features --- by a metric suited to the skew: ten
anomalies in $2010$ examples is a prevalence of $0.5\%$, so accuracy is useless
(\cref{sec:skew}) and precision, recall and $F_1$ are the measures.

\begin{workedexample}{Choosing $\varepsilon$ on the validation set}
\begin{center}
\small
\renewcommand{\arraystretch}{1.15}
\begin{tabular}{@{}rrrrrrr@{}}
\toprule
$\varepsilon$ & TP & FP & FN & precision & recall & $F_1$\\
\midrule
$10^{-3}$          & $10$ & $156$ & $0$ & $0.060$ & $1.0$ & $0.114$\\
$10^{-4}$          & $10$ & $16$  & $0$ & $0.385$ & $1.0$ & $0.556$\\
$3\times10^{-5}$   & $9$  & $8$   & $1$ & $0.529$ & $0.9$ & $0.667$\\
$10^{-5}$          & $7$  & $2$   & $3$ & $0.778$ & $0.7$ & $\mathbf{0.737}$\\
$3\times10^{-6}$   & $3$  & $0$   & $7$ & $1.000$ & $0.3$ & $0.462$\\
\bottomrule
\end{tabular}
\end{center}

$F_1$ is largest at $\varepsilon = 10^{-5}$. Applied once to the test set, that
threshold gives TP $= 8$, FP $= 0$, FN $= 2$: precision $1.0$, recall $0.8$,
$F_1 = 0.889$. It flags $0.07\%$ of the training engines.
\end{workedexample}

\begin{workedexample}{The same choice by hand, on twelve engines}
A real validation set is often small. Take nine good and three faulty validation
engines and sort them by $p(\vx)$:

\begin{center}
\small
\renewcommand{\arraystretch}{1.1}
\begin{tabular}{@{}lcr@{}}
\toprule
engine & faulty? & $p(\vx)$\\
\midrule
$(609.5,\ 3.3)$ & yes & $3.1\times10^{-7}$\\
$(411.0,\ 4.9)$ & yes & $4.1\times10^{-6}$\\
$(417.4,\ 5.1)$ & yes & $5.8\times10^{-6}$\\
$(470.4,\ 3.4)$ & no  & $3.2\times10^{-3}$\\
$(512.1,\ 3.4)$ & no  & $5.5\times10^{-3}$\\
\multicolumn{3}{@{}l}{\quad\dots\ seven more good engines, up to $1.2\times10^{-2}$}\\
\bottomrule
\end{tabular}
\end{center}

Only thresholds between consecutive values change the decisions. Any $\varepsilon$
between $5.8\times10^{-6}$ and $3.2\times10^{-3}$ flags exactly the three faulty
engines, $F_1 = 1$; lower thresholds miss some ($F_1 = 0.5$ or $0.8$), higher
ones add false alarms ($F_1 = 0.857$, falling to $0.429$). The geometric midpoint
of the perfect range is $1.4\times10^{-4}$ --- ten times the $\varepsilon$ chosen
on the full validation set. Twelve engines can say that some $\varepsilon$ in a
range spanning nearly three orders of magnitude is perfect; they cannot say
where in the range to put it.
\end{workedexample}

\begin{pitfall}{Very few labelled anomalies}
With only two or three known anomalies, use them all for validation and skip the
test set --- but then the reported $F_1$ is optimistic, because it was maximised
on the same examples it describes. Say so when reporting it. A few anomalies
slipping into the \emph{training} set, by contrast, do little harm: they move
$\hat\mu$ and $\hat\sigma^2$ only slightly.
\end{pitfall}

\begin{researchspotlight}{Chandola, Banerjee \& Kumar (2009) --- Anomaly detection: a survey}
\textbf{The problem.} Anomaly detection had grown up separately in statistics,
machine learning, data mining and a dozen application areas, each with its own
vocabulary, assumptions and methods.

\textbf{The idea.} Organise the field by the assumptions methods make about
normal data: statistical (normal data comes from a probability model, as in this
chapter), nearest-neighbour (normal data is dense), clustering-based (normal data
belongs to clusters), classification-based (labels exist for normal behaviour),
and others --- and by the type of anomaly: point, contextual, or collective.

\textbf{The result.} A widely cited map of the field, with the strengths,
weaknesses and computational cost of each family, and a catalogue of
applications from intrusion detection to medical monitoring.

\textbf{Why it is here.} The Gaussian detector of this chapter is the simplest
statistical method; the survey is the place to go when its assumptions fail ---
for instance when an anomaly is a \emph{sequence} of individually normal events.
\parencite{chandola2009}
\end{researchspotlight}

\section{Anomaly Detection or Supervised Learning?}
\label{sec:ad-vs-sl}

If some labelled anomalies exist, why not train a classifier on them?

\begin{center}
\small
\renewcommand{\arraystretch}{1.3}
\begin{tabular}{@{}L{5.3cm}L{5.3cm}@{}}
\toprule
\textbf{anomaly detection} & \textbf{supervised learning}\\
\midrule
very few positive examples ($0$--$20$), many negatives
  & enough positives and negatives to learn both\\
many \emph{different kinds} of anomaly; future ones may look like none seen
  & future positives resemble the training positives\\
fraud (attackers keep inventing), new engine faults, failing servers
  & spam, known defect types such as scratches, weather, known diseases\\
\bottomrule
\end{tabular}
\end{center}

\begin{keyinsight}{Model normal, or model the positives}
A classifier learns what the positive class looks like, and so can only find
positives resembling those it has seen. An anomaly detector learns what
\emph{normal} looks like, and flags anything else --- including kinds of anomaly
nobody has seen before. Many production systems run both: a classifier for the
known fraud patterns, and a detector for the unknown ones.
\end{keyinsight}

\begin{pitfall}{``Few positives'' is not the whole test}
A team with forty labelled defects reaches for anomaly detection because forty
is small. But if every future defect is a scratch, like the forty, a supervised
model --- perhaps with the data augmentation of \cref{sec:adding-data} --- will
find scratches better than any model of normality. Ask whether future anomalies
will resemble past ones before asking how many past ones you have.
\end{pitfall}

\section{Choosing and Transforming Features}
\label{sec:ad-features}

Feature choice matters more here than in supervised learning, because there is no
label to tell the model which features matter: a feature that does not separate
anomalies from normal data only adds noise to $p(\vx)$.

\subsection{Make features look Gaussian}

The model treats each feature as Gaussian. A feature with a long right tail ---
transaction amounts, response times, counts --- makes large but perfectly normal
values look improbable, and squeezes the real outliers together.

\begin{figure}[htbp]
  \centering
  \begin{minipage}[t]{0.48\linewidth}
    \centering
    \begin{tikzpicture}
      \begin{axis}[mlaxis, ybar, bar width=7pt, width=\linewidth, height=4.2cm,
        xmin=-0.25, xmax=6.25, ymin=0, title={raw $x$: skewness $3.33$}, ytick=\empty]
        \addplot[fill=MLBlue!55, draw=MLBlue] coordinates {(0.25,176)(0.75,310)(1.25,181)(1.75,139)(2.25,60)(2.75,41)(3.25,29)(3.75,22)(4.25,15)(4.75,4)(5.25,6)(5.75,3)};
      \end{axis}
    \end{tikzpicture}
  \end{minipage}\hfill
  \begin{minipage}[t]{0.48\linewidth}
    \centering
    \begin{tikzpicture}
      \begin{axis}[mlaxis, ybar, bar width=7pt, width=\linewidth, height=4.2cm,
        xmin=-3.25, xmax=3.25, ymin=0, title={$\ln x$: skewness $-0.08$}, ytick=\empty]
        \addplot[fill=MLTeal!55, draw=MLTeal] coordinates {(-2.75,1)(-2.25,8)(-1.75,21)(-1.25,63)(-0.75,159)(-0.25,234)(0.25,227)(0.75,171)(1.25,89)(1.75,19)(2.25,6)(2.75,2)};
      \end{axis}
    \end{tikzpicture}
  \end{minipage}
  \caption{$1000$ readings of a right-skewed feature, before and after taking
  logarithms. The raw histogram (which omits $14$ values above $6$) has a long
  tail; the logarithm is close to a symmetric bell, and a Gaussian model of it is
  reasonable.}
  \label{fig:skew-hist}
\end{figure}

Try $\ln x$, $\ln(x+c)$, $\sqrt{x}$ or $x^{1/3}$, judge by the histogram (or a
skewness statistic), and apply the \emph{same} transformation to the training,
validation and test data.

\subsection{Invent the feature that exposes the miss}

\term{Error analysis} for anomaly detection looks at the anomalies the detector
\emph{misses} and asks what makes them unusual.

\begin{workedexample}{A stuck process in a data centre}
For $500$ machines, CPU load and network traffic are strongly correlated
(correlation $0.97$): busy machines serve more requests. One machine shows
$85\%$ CPU with traffic of only $250$ --- a process stuck in a loop.

Each feature on its own is unremarkable: $1.52$ and $-1.37$ standard deviations
from the mean. The independent model gives it $p(\vx) = 4.5\times10^{-6}$, and
$79$ of the $500$ normal machines score \emph{lower}. No $\varepsilon$ flags the
stuck machine without also flagging those $79$.

Add the feature $x_3 = \text{CPU load}/\text{traffic}$. Normal machines have
$x_3 = 0.1006 \pm 0.0086$; the stuck one has $0.34$, which is $28$ standard
deviations away. Now it is the most improbable machine by an enormous margin.
\end{workedexample}

\section{The Multivariate Gaussian}
\label{sec:multivariate}

The ratio feature worked because someone understood the system well enough to
invent it. The multivariate Gaussian captures correlations between features
automatically.

\begin{defbox}{Multivariate Gaussian}
For $\vx \in \Real{n}$ with mean vector $\boldsymbol\mu$ and covariance matrix
$\Sigma$ ($n\times n$, symmetric, positive definite),
\[
  p(\vx;\boldsymbol\mu,\Sigma) = \frac{1}{(2\pi)^{n/2}|\Sigma|^{1/2}}
    \exp\!\Bigl(-\tfrac12(\vx - \boldsymbol\mu)^{\mathsf T}\Sigma^{-1}(\vx - \boldsymbol\mu)\Bigr).
\]
The maximum-likelihood estimates are $\hat{\boldsymbol\mu} = \frac1m\sum_i\xii{i}$
and $\hat\Sigma = \frac1m\sum_i(\xii{i} - \hat{\boldsymbol\mu})(\xii{i} -
\hat{\boldsymbol\mu})^{\mathsf T}$. The quadratic form in the exponent is the
squared \term{Mahalanobis distance}: distance measured in units of the data's
own spread, along its own axes.
\end{defbox}

\begin{keyinsight}{The independent model is a special case}
If $\Sigma$ is diagonal, with entries $\sigma_1^2,\dots,\sigma_n^2$, the
multivariate density factorises into exactly the product of \cref{sec:ad-alg}.
The independent model is a multivariate Gaussian whose ellipses are forced to
align with the axes; the full model lets them tilt.
\end{keyinsight}

\begin{figure}[htbp]
  \centering
  \begin{tikzpicture}
    \begin{axis}[mlaxis, width=0.84\linewidth, height=6.0cm, view={0}{90},
      xlabel={CPU load (\%)}, ylabel={network traffic},
      xmin=0, xmax=110, ymin=-200, ymax=1300, scaled y ticks=false,
      legend style={font=\scriptsize, at={(0.02,0.98)}, anchor=north west}]
      \addplot[only marks, mark=*, mark size=1pt, MLBlue!60, forget plot] coordinates {(76.4,673.5)(23.8,276.0)(23.4,248.2)(50.4,500.0)(47.5,483.4)(58.9,541.3)(67.5,686.7)(81.1,703.3)(55.2,517.6)(76.4,780.2)(51.4,512.9)(76.0,824.9)(21.0,229.4)(45.8,389.5)(58.9,569.8)(67.9,768.1)(27.5,301.9)(75.9,764.5)(46.6,489.0)(89.4,986.1)(69.9,663.4)(53.6,534.2)(20.1,204.8)(84.7,887.7)(77.6,708.4)(80.1,864.7)(38.1,384.1)(80.5,821.2)(20.2,189.9)(82.3,748.7)(80.1,857.0)(61.6,624.6)(68.8,706.0)(35.0,328.6)(62.1,721.5)(49.2,509.6)(46.3,436.5)(36.1,341.1)(57.4,483.3)(23.0,243.9)(21.4,186.8)(71.1,810.7)(63.1,667.5)(33.4,379.6)(40.6,338.8)(21.2,197.2)(68.3,663.4)(67.5,643.9)(39.2,373.5)(71.3,662.3)(26.5,222.0)(20.0,210.7)(62.1,625.4)(62.8,583.1)(47.7,494.0)(55.5,508.3)(53.8,506.7)(81.6,824.3)(24.2,246.9)(57.9,534.4)(84.3,835.9)(27.9,288.3)(24.7,252.4)(53.6,521.7)(25.0,252.1)(71.9,703.5)(49.8,570.9)(43.5,423.7)(86.3,723.6)(59.7,600.5)(82.1,832.6)(74.8,770.7)(29.1,269.3)(77.2,723.2)(90.0,1024.8)(66.0,605.7)(24.9,263.0)(37.1,371.6)(24.3,235.9)(85.2,916.5)(34.3,379.6)(60.1,619.0)(61.7,606.8)(54.3,555.3)(56.1,578.5)(37.8,367.6)(68.9,738.5)(43.3,423.6)(68.5,613.7)(25.4,263.3)(37.9,371.6)(22.7,234.3)(75.4,659.8)(24.6,293.6)(71.2,658.1)(36.7,349.2)(34.8,347.8)(21.5,209.0)(74.5,663.3)(65.5,674.4)(28.7,292.8)(32.0,278.9)(83.7,953.0)(49.7,479.8)(35.9,347.1)(85.8,802.4)(75.1,734.1)(22.6,200.2)(59.6,628.1)(20.8,189.9)(39.7,392.2)(36.4,387.1)(79.0,690.9)(56.4,570.3)(71.0,690.6)(20.4,218.7)(77.9,843.9)(64.1,555.0)(24.4,231.7)(82.5,917.4)(57.7,617.6)(80.7,875.5)(86.0,909.2)(43.7,406.2)(36.8,373.7)};
      \costcontour{contour/draw color=MLAmber, very thick, dashed, forget plot, domain=0:110, y domain=-200:1300}{4.1985}{((x-54.0121)/20.3907)^2+((y-540.6261)/211.4579)^2}
      \addlegendimage{MLAmber, very thick, dashed}\addlegendentry{independent model, through the stuck machine}
      \costcontour{contour/draw color=MLTeal, very thick, forget plot, domain=0:110, y domain=-200:1300}{22.9960}{0.04238942*(x-54.0121)^2+2*(-0.00396992)*(x-54.0121)*(y-540.6261)+0.0003941609*(y-540.6261)^2}
      \addlegendimage{MLTeal, very thick}\addlegendentry{multivariate model, least likely normal machine}
      \addplot[only marks, mark=x, mark size=3pt, MLRose, very thick, forget plot] coordinates {(85,250)};
      \node[font=\tiny, MLRose, anchor=west] at (axis cs:87,250) {stuck};
    \end{axis}
  \end{tikzpicture}
  \caption{The data-centre machines (dots) and the stuck one (cross). Dashed: the
  independent model's contour through the stuck machine --- an axis-aligned
  ellipse with $79$ normal machines outside it, all scored as more anomalous.
  Solid: the multivariate model's contour through the least likely
  \emph{normal} machine, tilted along the correlation. The stuck machine lies
  far outside it.}
  \label{fig:multivariate}
\end{figure}

On the data-centre machines the estimated correlation is $0.9712$. The
multivariate model gives the stuck machine $p = 3.9\times10^{-36}$, a squared
Mahalanobis distance of $145.5$, while the least likely normal machine has
$p = 1.6\times10^{-9}$ and squared distance $23.0$. The anomaly is found from the
raw features, with no invented ratio.

\begin{center}
\small
\renewcommand{\arraystretch}{1.3}
\begin{tabular}{@{}L{3.1cm}L{4.1cm}L{4.1cm}@{}}
\toprule
 & independent Gaussians & multivariate Gaussian\\
\midrule
correlations & missed unless encoded as features & captured automatically\\
parameters & $2n$ ($200$ for $n = 100$) & $n + n(n+1)/2$ ($5150$ for $n = 100$)\\
cost & linear in $n$ & needs $\Sigma^{-1}$: about $n^3$\\
data needed & works with $m$ small & needs $m > n$, ideally $m \gg n$\\
\bottomrule
\end{tabular}
\end{center}

\begin{pitfall}{A singular covariance matrix}
If $m \le n$, or if some features are duplicates or exact linear combinations of
others, $\hat\Sigma$ is singular and $\Sigma^{-1}$ does not exist --- the
singular-matrix problem of \cref{sec:singular} again. Remove redundant features,
or add a small multiple of the identity to $\hat\Sigma$, the same regularization
idea as ridge regression in \cref{ch:ch07}.
\end{pitfall}

\begin{labbox}{Lab 14 --- Anomaly detection}
Implement \code{estimate\_gaussian} returning $\hat\mu$ and $\hat\sigma^2$ per
feature, and a vectorised \code{p(X)}. Reproduce the estimates and the three
engine scores of \cref{sec:ad-alg}. Implement \code{select\_threshold} that scans
candidate $\varepsilon$ values on the validation set and returns the best $F_1$,
and reproduce $\varepsilon = 10^{-5}$ with $F_1 = 0.737$.

Then fit the multivariate Gaussian to the data-centre machines, plot both models'
contours as in \cref{fig:multivariate}, and confirm that only the multivariate
model isolates the stuck machine. Finally apply the detector to a dataset with
many features, where the provided labels let you report precision and recall.
\end{labbox}

\section*{Chapter Summary}
\addcontentsline{toc}{section}{Chapter Summary}

\begin{itemize}
  \item Anomaly detection models normal data with a density $p(\vx)$ and flags
        examples with $p(\vx) < \varepsilon$.
  \item The maximum-likelihood estimates of a Gaussian's parameters are the
        sample mean and the sample variance (dividing by $m$).
  \item The basic detector multiplies one Gaussian per feature; its boundaries are
        axis-aligned ellipses. On E2, an engine $3.6$ and $3.2$ standard
        deviations out scores $1.0\times10^{-7}$ and is flagged.
  \item Train on unlabelled, mostly normal data; choose $\varepsilon$ by $F_1$ on a
        few labelled anomalies. On E2, $\varepsilon = 10^{-5}$ gives validation
        $F_1 = 0.737$ and test $F_1 = 0.889$. Tiny validation sets leave
        $\varepsilon$ poorly determined.
  \item Prefer anomaly detection when positives are rare \emph{and} varied, so
        future ones may be new; prefer supervised learning when they are
        plentiful or alike.
  \item Transform skewed features towards a bell shape, and build features from
        the anomalies the detector misses --- a ratio that puts a stuck machine
        $28$ standard deviations out.
  \item The multivariate Gaussian captures correlations; the independent model is
        its diagonal special case. It isolates the stuck machine from the raw
        features, at the cost of $O(n^2)$ parameters and a covariance matrix that
        must be invertible.
\end{itemize}

\section*{Exercises}
\addcontentsline{toc}{section}{Exercises}

\exercisesection{Short Questions}

\sq Explain anomaly detection by density estimation and give three applications.

\sq Write the Gaussian density and say what $\mu$ and $\sigma$ control.

\sq Why can a density exceed $1$, and why does that not matter here?

\sq What assumption does multiplying the per-feature densities make?

\sq Why is accuracy unsuitable for choosing $\varepsilon$?

\sq Why do a few anomalies in the training set do little harm?

\sq When should supervised learning be preferred to anomaly detection?

\sq Why should features be transformed towards a Gaussian shape, and which
transformations are common?

\sq How does error analysis lead to new features in anomaly detection?

\sq How many parameters does a multivariate Gaussian have in $n$ dimensions?

\exercisesection{Analytical Problems}

\ap Derive the maximum-likelihood estimates of $\mu$ and $\sigma^2$ for a
Gaussian, and show that $\hat\sigma^2$ is biased by a factor $(m-1)/m$.

\ap Compute $\hat\mu$ and $\hat\sigma^2$ for $x = \{2, 4, 4, 4, 5, 5, 7, 9\}$, and
the density at $x = 10$.

\ap Using E2's estimates, compute $p(\vx)$ for $\vx = (450,\ 3.0)$ and decide
whether it is an anomaly at $\varepsilon = 10^{-5}$.

\ap Show that the boundary $p(\vx) = \varepsilon$ of the independent model in two
dimensions is an ellipse, and find its semi-axes for E2 at $\varepsilon =
10^{-5}$.

\ap Show that a multivariate Gaussian with diagonal covariance factorises into a
product of one-dimensional Gaussians.

\ap Explain, using the Mahalanobis distance, why the stuck machine is far from
the data under the multivariate model but not under the independent one.

\exercisesection{Design Problems}

\dpr Design a fraud detector for a mobile payment service: the features
(including at least one ratio), the training data, how you obtain labelled frauds
for validation, and how you choose $\varepsilon$ given the cost of blocking a
genuine payment.

\dpr Design an anomaly detector for sewing machines in a factory using vibration
and power sensors, including how you would detect a machine that degrades slowly
over weeks.

\dpr You have $40$ labelled defects and $50{,}000$ good circuit boards. Design an
experiment to decide between anomaly detection and a supervised classifier.

\dpr Design an automatic procedure that chooses a transformation for each
feature before fitting $p(\vx)$, and a test that it improves validation $F_1$.

\exercisesection{Programming Exercises}

\pe Implement the independent-Gaussian detector and reproduce the
$\varepsilon$ table of \cref{sec:ad-eval}.

\pe Implement \code{select\_threshold}, scanning candidate thresholds between the
sorted validation scores, and reproduce the twelve-engine example.

\pe Fit independent and multivariate Gaussians to the data-centre machines and
reproduce \cref{fig:multivariate}.

\pe Generate a skewed feature, fit the detector with and without a log transform,
and compare their validation $F_1$ on injected anomalies.

\exercisesection{Self-Test}

\qz The maximum-likelihood estimate of $\sigma^2$ divides by: (a)~$m-1$;
(b)~$m$; (c)~$n$; (d)~$\sqrt m$.

\qz An example is flagged as anomalous when: (a)~$p(\vx) > \varepsilon$;
(b)~$p(\vx) < \varepsilon$; (c)~$p(\vx) = 1$; (d)~its label is $1$.

\qz $\varepsilon$ is chosen on: (a)~the training set; (b)~the validation set,
by $F_1$; (c)~the test set; (d)~no data at all.

\qz Anomaly detection is preferred over supervised learning when: (a)~positives
are plentiful and alike; (b)~positives are rare and varied; (c)~the data is
labelled; (d)~accuracy is high.

\qz The independent-Gaussian model is the multivariate Gaussian with:
(a)~a diagonal covariance; (b)~an identity mean; (c)~zero variance; (d)~a singular
covariance.

\qz The multivariate Gaussian with $n = 100$ features has about: (a)~$200$
parameters; (b)~$5150$; (c)~$100$; (d)~$10^4$.

\qz A long right tail in a feature is best handled by: (a)~deleting the feature;
(b)~a log or similar transform; (c)~squaring it; (d)~nothing.

\furtherreading

\textcite{james2021} §4.4.1--4.4.2 is the primary reading for the probability
models of this chapter: it develops the one-dimensional Gaussian density and its
parameter estimates, then the multivariate Gaussian with its covariance matrix,
in the setting of linear discriminant analysis, where the same densities are used
to classify rather than to flag outliers. \textcite{bishop2006} §2.3 is the
supplementary reading: it treats the multivariate Gaussian in depth --- its
geometry, its Mahalanobis distance and elliptical contours, and in §2.3.4 the
maximum-likelihood estimates of its mean and covariance --- and is the right place
to see why the independent model of \cref{sec:ad-alg} is its diagonal special
case.

\bigskip
\noindent\emph{Next:} \Cref{ch:ch15} opens Part~VI with recommender systems,
where a matrix of ratings with most entries missing is completed from the few
that are known.
